% !TeX program = pdflatex
%
% Documentation of the neutronics homogenisation subsystem
% (homogenise.py, homogenise_assembly.py, component_material_zones.py)
% of the ZeroLevelProduct CAD assembly toolchain.
%
% Written 2026-07-22. Companion to cost_reporting.tex, which documents the
% other consumer of the same component specification list.
%
% extarticle (from the extsizes bundle) is a drop-in clone of article that
% accepts sizes below 10pt. Every other size in this file is RELATIVE
% (\footnotesize, \scriptsize, ...), so the base size set here scales the whole
% document -- body, tables, code listings and diagram labels alike.
% If extsizes is not installed, revert to:  \documentclass[10pt,a4paper]{article}
\documentclass[9pt,a4paper]{extarticle}

\usepackage[T1]{fontenc}
\usepackage[utf8]{inputenc}
\usepackage[margin=2.5cm]{geometry}
\usepackage{amsmath}
\usepackage{amssymb}
\usepackage{booktabs}
\usepackage{longtable}
\usepackage{array}
\usepackage{xcolor}
\usepackage{listings}
\usepackage{tikz}
\usetikzlibrary{arrows.meta,positioning,fit,backgrounds}
\usepackage[hidelinks]{hyperref}

\definecolor{codebg}{gray}{0.96}
\definecolor{codekw}{RGB}{0,80,160}
\definecolor{codecm}{RGB}{100,120,100}
\definecolor{codestr}{RGB}{160,60,60}
\definecolor{warnbg}{RGB}{255,247,230}
\definecolor{warnrule}{RGB}{200,140,40}

\lstset{
  basicstyle=\ttfamily\scriptsize,
  backgroundcolor=\color{codebg},
  keywordstyle=\color{codekw}\bfseries,
  commentstyle=\color{codecm}\itshape,
  stringstyle=\color{codestr},
  showstringspaces=false,
  breaklines=true,
  frame=single,
  rulecolor=\color{gray!40},
  framesep=4pt,
  xleftmargin=6pt,
  xrightmargin=2pt,
  literate={³}{{\textsuperscript{3}}}1 {·}{{$\cdot$}}1 {→}{{$\rightarrow$}}1
           {ρ}{{$\rho$}}1 {Σ}{{$\Sigma$}}1 {—}{{---}}1 {’}{{'}}1,
}
\lstdefinestyle{py}{language=Python}

\newcommand{\code}[1]{\texttt{\small #1}}
\newcommand{\file}[1]{\texttt{\small #1}}

% A boxed caveat, for the findings that a reader must not miss.
% Plain LaTeX (lrbox + minipage + fcolorbox) rather than a tikz node, so that it
% survives page breaks in the surrounding text and needs no extra package.
\newsavebox{\caveatbox}
\newenvironment{caveat}[1]%
  {\par\medskip\noindent
   \begin{lrbox}{\caveatbox}%
   \begin{minipage}{\dimexpr\linewidth-2\fboxsep-2\fboxrule\relax}%
   \textbf{#1}\par\medskip}%
  {\end{minipage}%
   \end{lrbox}%
   \fcolorbox{warnrule}{warnbg}{\usebox{\caveatbox}}\par\medskip}

\title{\textbf{Neutronics Homogenisation in the Zero-Level Product Toolchain}\\[0.4em]
       \large Atom-conserving equivalent cylinders from CAD geometry}
\author{Documentation of \file{homogenise.py}, \file{homogenise\_assembly.py}\\
        and \file{component\_material\_zones.py}}
\date{\today}

\begin{document}
\maketitle

\begin{abstract}
\noindent
This document describes the neutronics homogenisation subsystem of the
zero-level reactor product model. The subsystem takes the \emph{same} component
specification list that drives the CAD assembly, measures the true CAD volume of
every material zone of every component, and replaces each component with an
equivalent right circular cylinder whose per-nuclide number densities conserve
the number of atoms of every nuclide present. The equivalent cylinders are
placed at the components' resolved world positions --- in particular at the same
radial position --- inside a sodium pool bounded by a natively modelled reactor
vessel and top plate. The result is a plain, JSON-serialisable model intended to
be handed to a Monte Carlo transport code for a $k$-eigenvalue, fluence or
activation calculation. This document covers the conservation identity and its
unit handling, the zone model, the four component treatments, the placement
convention, the structure of the output, the visualisation path, and the
verification strategy --- including a precise statement of which checks can fail
and which are true by construction. A worked example on the ESFR-SMR reference
assembly is included, together with an audit of the known limitations of the
current model: three components whose equivalent cylinder does not enclose them,
a top plate that is over-carved by the equivalent cylinders, a fuel composition
applied to roughly four times the height it was derived for, and the absence
of both an outer boundary and a transport-code exporter.
\end{abstract}

\tableofcontents
\medskip

% =====================================================================
\section{Purpose and design intent}
% =====================================================================

The CAD side of this toolchain builds a geometrically faithful reactor assembly.
That geometry is not what a Monte Carlo transport calculation should consume: an
intermediate heat exchanger with $112$ individually modelled tubes is expensive
to track and contributes detail that a first-pass $k$-eigenvalue or fluence
calculation does not need.

The homogenisation subsystem produces the \emph{simplified} model instead. Its
governing principle is stated once and enforced everywhere:

\begin{quote}
Every component is replaced by an equivalent cylinder containing exactly the
same number of atoms of every nuclide, positioned where the real component sits.
\end{quote}

Two consequences follow, and they are the reason the subsystem is worth having
rather than hand-writing the simplified model:

\begin{itemize}
  \item \textbf{The simplified model cannot drift from the CAD model.} Both are
    generated from one specification list. Change a wall thickness and both
    change together.
  \item \textbf{Material inventory is exact, not approximate.} Volumes are
    \emph{measured} from the built CAD solids rather than estimated from
    nominal dimensions, so the steel mass in the neutronics model is the steel
    mass in the drawing.
\end{itemize}

What is deliberately \emph{not} preserved is shape. A component's material is
redistributed uniformly through its equivalent cylinder; where the component
does not fill the cylinder, the remainder is filled with a declared coolant
rather than left as void. Section~\ref{sec:audit} quantifies where this
redistribution is largest.

% =====================================================================
\section{Where it sits in the pipeline}
% =====================================================================

\subsection{Two entry points, one specification list}

The specification list --- a plain Python list of dictionaries, one per
component --- feeds two independent consumers:

\begin{center}
\begin{tabular}{@{}ll@{}}
\toprule
\code{assemble\_objects(SPECS)} & the CAD assembly (and, optionally, the cost report) \\
\code{homogenise\_assembly(SPECS)} & the homogenised neutronics model \\
\bottomrule
\end{tabular}
\end{center}

\noindent
They are siblings: neither calls the other, and each ignores the keys the other
needs. \code{assemble\_objects} ignores \code{"materials"},
\code{"neutronics\_treatment"} and \code{"hom\_cylinder"};
\code{homogenise\_assembly} ignores nothing it is given but adds no geometry.
The only additions a user makes to an existing CAD specification list in order
to homogenise it are material assignments (Section~\ref{sec:specs}).

Both entry points internally run the \emph{same} two pre-passes --- the
component resolver, which converts \code{at\_radius} / \code{at\_angle\_deg}
into Cartesian \code{center\_coords}, and \code{auto\_generate\_topplate\_holes},
which punches the top-plate penetrations. Running them in the homogeniser as
well is what guarantees a component lands at the same world position in both
models, and that the top plate is measured as a perforated slab rather than a
solid one.

\begin{caveat}{Order independence}
\code{assemble\_objects} does not deep-copy the caller's dictionaries: it
resolves them and adds top-plate holes \emph{in place}. Calling it before
\code{homogenise\_assembly} on the same list therefore hands the homogeniser
already-resolved dictionaries. This is safe: the hole-merge rule
(``a user hole claiming a component suppresses the auto hole'') makes the
pre-pass idempotent, and the two call orders were verified to produce
bit-identical plate and pool volumes (Section~\ref{sec:repro}).
\end{caveat}

\subsection{Data flow}

\begin{figure}[h]
\centering
\begin{tikzpicture}[
  node distance=6mm,
  box/.style={draw, rounded corners=2pt, align=center, font=\scriptsize,
              minimum height=8mm, inner sep=4pt, fill=white},
  src/.style={box, fill=blue!6},
  out/.style={box, fill=green!8},
  arr/.style={-{Latex[length=2mm]}, gray!70, thick}
]
\node[src] (specs) {\textbf{SPECS}\\component dicts\\(geometry + materials)};
\node[box, right=13mm of specs] (res) {\code{resolve()}\\
      \code{auto\_generate\_}\\\code{topplate\_holes()}};
\node[box, right=13mm of res] (zones) {\code{zone\_model(spec)}\\
      \code{MATERIAL\_ZONES} $\rightarrow$ zones\\
      + envelope; else generic};

\node[src, below=12mm of specs] (matlib) {\file{materials.py}\\\code{MATERIALS}};
\node[box, right=13mm of matlib] (nd) {\code{material\_number\_}\\\code{densities()}\\
      $n_{m,i}$ in atoms/b$\cdot$cm};
\node[box, right=13mm of nd] (hom) {\code{homogenise\_volumes()}\\
      $n_{\mathrm{hom},i}=\sum_m \frac{V_m}{V_{\mathrm{cyl}}} n_{m,i}$};

\node[box, below=12mm of matlib] (place) {\code{\_world\_cylinder()}\\
      local $+$ \code{center\_coords}};
\node[out, right=13mm of place] (model) {\textbf{model} (dict)\\
      components, pool,\\vessel, plate, settings, qa};
\node[out, right=13mm of model] (cons) {\code{build\_cad()} $\rightarrow$ CAD\\
      \code{out\_path=} $\rightarrow$ JSON\\
      \code{print\_report()}};

\draw[arr] (specs) -- (res);
\draw[arr] (res) -- (zones);
\draw[arr] (matlib) -- (nd);
\draw[arr] (nd) -- (hom);
\draw[arr] (zones.south) -- ++(0,-4mm) -| (hom.north);
\draw[arr] (hom.south) -- ++(0,-4mm) -| (model.north);
\draw[arr] (place) -- (model);
\draw[arr] (model) -- (cons);
\end{tikzpicture}
\caption{Data flow of the homogeniser. Blue boxes are inputs the user supplies
or configures; green boxes are produced artefacts. The \code{model} dictionary
is a neutral intermediate: it contains no CadQuery objects and is fully
JSON-serialisable, so every downstream consumer --- including the future
transport-code exporter --- reads it rather than recomputing.}
\end{figure}

% =====================================================================
\section{Quick start}
% =====================================================================

Given a specification list \code{SPECS} in which every component carries a
material assignment, the complete workflow is:

\begin{lstlisting}[style=py]
from homogenise_assembly import (homogenise_assembly, build_cad,
                                 check_against_cad, print_report)
from ocp_vscode import show

model = homogenise_assembly(SPECS)                       # 1. compute
print_report(model, check_against_cad(model, SPECS))     # 2. inspect  (optional)
show(build_cad(model))                                   # 3. visualise (optional)
\end{lstlisting}

\noindent
Note what is \emph{not} imported. The material library is the default value of
the \code{materials} argument, so the user never handles it as an object ---
materials are named by string inside the specification dictionaries, exactly as
\code{obj\_type} is. Pass \code{materials=} only to override the library, for
instance with a temperature-perturbed composition set. The pool material
likewise defaults to \code{"na\_primary"}.

To keep the model on disk, and to export the simplified geometry as STEP:

\begin{lstlisting}[style=py]
model = homogenise_assembly(SPECS, out_path="output/homogenised_model.json")
build_cad(model, export_path="output/homogenised.step", units="m")
\end{lstlisting}

\noindent
\code{units} is required for the STEP export and \emph{only} for it, because the
file format demands a length declaration. The homogenisation itself never needed
one (Section~\ref{sec:units}).

To view the real and simplified geometries superimposed --- the most direct way
to judge whether the simplification is acceptable:

\begin{lstlisting}[style=py]
show(assemble_objects(SPECS),
     build_cad(homogenise_assembly(SPECS)),
     names=["heterogeneous", "homogenised"],
     alphas=[1.0, 0.35])
\end{lstlisting}

% =====================================================================
\section{What a specification must carry}
\label{sec:specs}
% =====================================================================

\subsection{Material assignment (required)}

Every component must bind each of its material zones to a named material. Two
forms are accepted:

\begin{lstlisting}[style=py]
# multi-zone component: one material per declared zone
"materials": {"structure": "ss316", "tube_side": "na_secondary",
              "shell_side": "na_primary", "_filler": "na_primary"}

# single-zone component: the shorthand
"material": "ss316", "filler": "na_primary"
\end{lstlisting}

\noindent
The shorthand is accepted \emph{only} when the component declares exactly one
zone; using it on a multi-zone component raises a \code{KeyError} naming the
zones that need binding. The zone names available per component type are:

\begin{center}
\begin{tabular}{@{}ll@{}}
\toprule
\code{obj\_type} & zones \\
\midrule
\code{ihx}          & \code{structure}, \code{tube\_side}, \code{shell\_side} \\
\code{diagrid}      & \code{shell}, \code{cavity} \\
\code{reactor\_core}& \code{core} \\
everything else     & \code{structure} (registered) or \code{solid} (generic) \\
\bottomrule
\end{tabular}
\end{center}

\noindent
These need not be memorised. A wrong or missing zone name raises a
\code{KeyError} that lists the component's actual zone names, and the same is
true in reverse for a material bound to a zone that does not exist. Material
names must exist in \code{MATERIALS} (\file{materials.py}); an unknown name
raises a \code{KeyError} identifying the offending material.

\subsection{The filler (the closure rule)}

An equivalent cylinder is almost always larger than the sum of the component's
declared zones. The remainder,
\[
V_{\mathrm{filler}} \;=\; V_{\mathrm{cyl}} - \sum_z V_z ,
\]
is not void. In a pool-type reactor it is the surrounding sodium. It is assigned
with the reserved key \code{"\_filler"} in the \code{materials} block, or with
the top-level \code{"filler"} key on a single-material component. If no filler
is given, the remainder is left as vacuum and reported in the result as
\code{unfilled}, so the omission is visible rather than silent.

A \emph{negative} remainder --- declared zones larger than the cylinder --- is
always an error and raises immediately, naming both volumes. It means the
envelope is too small for the zones the component declared, and is repaired
either by enlarging the envelope or by fixing the zone volumes.

\subsection{Optional keys}

\begin{center}
\begin{tabular}{@{}lp{10.5cm}@{}}
\toprule
key & effect \\
\midrule
\code{"neutronics\_treatment"} &
  Forces \code{"cylinder"}, \code{"smear"} or \code{"skip"}, overriding the
  default for that \code{obj\_type} (Section~\ref{sec:treatments}). \\
\code{"hom\_cylinder"} &
  Partial envelope override, e.g. \code{\{"radius": 2.2\}}. Merged over the
  envelope the zone model declared. Used to shrink a cylinder clear of a
  neighbour; atoms remain conserved because only $V_{\mathrm{cyl}}$ changes and
  the filler absorbs the difference. \\
\bottomrule
\end{tabular}
\end{center}

\noindent
Both are stripped before the specification reaches \code{build\_solid}, along
with the material keys, by \code{\_NON\_GEOMETRY\_KEYS} in \file{homogenise.py}
--- otherwise the profile-based build path would reject them as unexpected
keyword arguments.

% =====================================================================
\section{The conservation identity}
\label{sec:units}
% =====================================================================

\subsection{Statement}

Let a component occupy volume $V_m$ of material $m$, each material having
per-nuclide number densities $n_{m,i}$ (atoms per barn-centimetre). The
equivalent cylinder of volume $V_{\mathrm{cyl}}$ is assigned

\begin{equation}
n_{\mathrm{hom},i} \;=\; \sum_m \frac{V_m}{V_{\mathrm{cyl}}}\, n_{m,i} .
\label{eq:hom}
\end{equation}

\noindent
Multiplying through by $V_{\mathrm{cyl}}$ gives the conservation statement
directly:

\begin{equation}
V_{\mathrm{cyl}}\, n_{\mathrm{hom},i} \;=\; \sum_m V_m\, n_{m,i}
\qquad \text{for every nuclide } i .
\label{eq:cons}
\end{equation}

\noindent
The left side is the atom content of the equivalent cylinder; the right side is
the atom content of the real component. They are equal by construction, for
every nuclide independently --- not merely in total, and not merely for the
dominant isotope.

\subsection{Why no unit is needed}

The ratio $V_m/V_{\mathrm{cyl}}$ in Equation~\ref{eq:hom} is dimensionless.
Scaling every length in the model by a factor $k$ scales every volume by $k^3$,
in numerator and denominator alike, leaving $n_{\mathrm{hom},i}$ unchanged.
The subsystem is therefore \emph{unit-agnostic} in exactly the sense
\code{assemble\_objects} is: the model may be in metres, millimetres or
furlongs, provided one unit is used consistently, and no number is ever
rescaled.

The resulting $n_{\mathrm{hom},i}$ emerges in the same units as the library's
$n_{m,i}$, namely atoms/b$\cdot$cm --- an absolute density that does not depend
on the model's length unit at all. A length unit is required in exactly two
places, both of them at the boundary with a format that demands one: the
\code{units} argument of \code{build\_cad} when writing STEP, and the future
transport-code export (Section~\ref{sec:openmc}).

This invariance is asserted directly in the module self-test, which repeats the
mixing calculation with every length scaled by $10^3$ and requires the
number-density drift to be below $10^{-12}$; it measures exactly $0$.

\subsection{Material definitions}

\code{material\_number\_densities()} converts a library entry into
$n_{m,i}$. Two forms are accepted:

\begin{lstlisting}[style=py]
{"kind": "number_density", "nuclides": {"Fe56": 5.9e-2, ...}}          # atoms/b-cm
{"kind": "mass_density", "density": 0.83,
 "nuclides": {"Na23": {"ao": 1.0, "M": 22.99}}}                        # g/cm3 + atom fractions
\end{lstlisting}

\noindent
For the mass-density form the atom fractions are normalised, the mean molar mass
$\bar{M}=\sum_i a_i M_i$ formed, and the total number density obtained as
$n_{\mathrm{tot}} = \rho N_A / \bar{M}$, converted to atoms/b$\cdot$cm by the
factor $10^{-24}$. Avogadro's number is defined once, in \file{homogenise.py},
and shared with the cost reporter so that both workflows use one constant.

% =====================================================================
\section{The zone model}
% =====================================================================

\subsection{What a zone is, and the contract}

A \emph{zone} is a named region of a component made of one material: a lump of
steel, a bore, a cavity, a pool side. A zone function takes the component
dictionary and returns

\begin{lstlisting}[style=py]
{"zones":    [Zone(name, role, volume), ...],              # model units cubed
 "cylinder": {"radius": R, "z_bottom": z0, "height": H},   # the envelope
 "solids":   {zone_name: Shape, ...}}                      # optional
\end{lstlisting}

\noindent
\code{role} is \code{"solid"} or \code{"fluid"} and is informational to the
homogeniser --- atom conservation treats the two identically. It exists because
the cost reporter bills only the solid zones, a bill of materials listing parts
to buy and not the coolant they sit in.

\begin{caveat}{Frame convention}
Everything a zone function returns is in the component's \textbf{local} frame
--- the frame \code{build\_solid()} works in \emph{before} it places the origin
at \code{center\_coords}. A zone function must never bake a world position into
its envelope. The assembly-level homogeniser performs the translation, and
doing it twice is precisely the class of error that
\code{check\_against\_cad()} exists to catch.
\end{caveat}

\subsection{Registered and generic components}

Registration in \code{MATERIAL\_ZONES} (\file{component\_material\_zones.py}) is
\emph{optional}. A component with no entry is homogenised generically by
\code{generic\_zones()}: it is built through the very same \code{build\_solid()}
call the assembler uses, treated as one zone named \code{solid}, and given its
own bounding cylinder as the envelope. That covers every single-material
component on all three build paths --- premade components, 2D profile plus
operation, and CadQuery 3D primitives --- with no registration required.

An entry is added only when a component has internal structure that one material
cannot express. The intermediate heat exchanger is the motivating case: its
bundle holds secondary sodium inside the tubes and primary sodium outside them,
within one fused solid.

\subsection{The bounding cylinder of a generic component}

\code{bounding\_cylinder()} deserves a note, because the obvious implementation
is wasteful. The bounding box's \emph{corner} distance
$\max\sqrt{x^2+y^2}$ always contains the solid, but for an axisymmetric
component it is a factor $\sqrt{2}$ too wide --- which costs nothing in atoms,
yet dilutes the homogenised densities over twice the volume they belong in. The
implementation therefore tries the tighter radius (the largest bounding-box
extent) first, and keeps it only if a single boolean cut confirms that nothing
protrudes:

\begin{lstlisting}[style=py]
if solid.cut(probe).Volume() <= tol * max(solid.Volume(), 1.0):
    radius = r_tight
\end{lstlisting}

\noindent
Components whose axis is not the local $z$ axis receive a generous envelope from
this and should be tightened with \code{spec["hom\_cylinder"]}.

\subsection{The registered components}

\begin{center}
\begin{tabular}{@{}llp{7.6cm}@{}}
\toprule
\code{obj\_type} & zones & envelope \\
\midrule
\code{ihx} & \code{structure}, \code{tube\_side}, \code{shell\_side} &
  from the plenum domes and bundle shell \\
\code{reactor\_core} & \code{core} &
  the specified \code{radius} and \code{height} \\
\code{diagrid} & \code{shell}, \code{cavity} &
  outer diameter, grown by the nozzle boss height when bosses are present \\
\code{primary\_pump} & \code{structure} &
  \code{barrel\_radius} $\times$ \code{barrel\_height} \\
\code{strongback} & \code{structure} &
  skirt outer radius and total height, or the profile's extent \\
\code{above\_core\_structure} & \code{structure} &
  covers the offset top cylinder and any CRDL / bottom-plate protrusion \\
\code{redan} & \code{structure} &
  the profile's extent (but the assembly smears it; see below) \\
\code{reactor\_vessel} & \code{structure} &
  supplies the exact CAD steel volume the native model must conserve \\
\code{reactor\_top\_plate} & \code{structure} &
  as above; must be measured \emph{after} the hole pre-pass \\
\bottomrule
\end{tabular}
\end{center}

\noindent
Two of these are subtler than they look. The diagrid's steel is built through
the assembly's own builder rather than \code{create\_diagrid()}, so that the
pump nozzle bosses and bores injected by the resolver are part of the measured
steel; and its wall-thickness precedence must match the builder's, an individual
\code{wall\_t\_side} winning over the uniform \code{wall\_t} shortcut. Reading
them the other way round measures a different shell than the one built. The
IHX's tube-side region is reconstructed here from the same parameters
\code{create\_ihx()} uses and then cut against the built structure, so the two
must be kept in step; the closure figure reported by \code{homogenise()} is the
guard, since a drifting reconstruction stops the zones summing to the envelope.

% =====================================================================
\section{The four treatments}
\label{sec:treatments}
% =====================================================================

Not every component should become a cylinder. The assembly-level homogeniser
dispatches on \code{obj\_type}, overridable per specification with
\code{"neutronics\_treatment"}:

\begin{center}
\begin{tabular}{@{}llp{8.2cm}@{}}
\toprule
treatment & default for & what happens \\
\midrule
\code{cylinder} & everything else &
  equivalent cylinder at the resolved world position \\
\code{vessel} & \code{reactor\_vessel} &
  kept native: cylindrical annulus plus an equivalent flat bottom conserving the
  CAD steel volume, heads included \\
\code{plate} & \code{reactor\_top\_plate} &
  kept native: slab minus the penetrating cylinders, with the CAD steel smeared
  through the remaining cell \\
\code{smear} & \code{redan} &
  steel dissolved uniformly into the pool \\
\code{skip} & --- &
  excluded from the model entirely \\
\bottomrule
\end{tabular}
\end{center}

\noindent
The redan is the instructive case. It is a thin revolved shell spanning the
whole pool; an equivalent cylinder of its outer radius would swallow the heat
exchangers and pumps whole. Smearing its steel into the pool conserves its atoms
without claiming a volume it does not exclusively occupy.

\subsection{Geometry caching}

Identical components are homogenised once. \code{\_cache\_key()} hashes the
specification with the placement keys stripped --- \code{obj\_id},
\code{at\_radius}, \code{center\_coords}, \code{rotation\_angles} and the rest
--- so three identically parameterised heat exchangers at three different
azimuths share one cached result. In the reference assembly this halves the
build: three IHXs and three pumps become two builds, reported in the log as
\code{reusing ... (identical geometry)}.

% =====================================================================
\section{Placement}
% =====================================================================

\code{build\_solid()} rotates a component about its \textbf{local origin} and
then translates that origin onto \code{center\_coords}. Hence, for a world point
$\mathbf{p}$ and local point $\mathbf{q}$,
\[
\mathbf{p} \;=\; R(\text{roll},\text{pitch},\text{yaw})\,\mathbf{q}
              \;+\; \mathbf{c},
\]
and for a $z$-aligned envelope centred on the component's own axis the yaw term
\emph{drops out entirely}: rotating about the local origin spins the component
about that very axis, and a circular envelope is invariant under it. The world
envelope is therefore the local envelope translated by \code{center\_coords},
which is what \code{\_world\_cylinder()} computes. Polar placement
(\code{center\_coords\_pol}) is converted to Cartesian first.

Roll or pitch tilt the component out of the $z$ direction, which no $z$-aligned
cylinder can represent. That case emits a warning and places the cylinder
upright.

\begin{caveat}{Radial position is preserved exactly}
Because placement is a pure translation by \code{center\_coords}, an equivalent
cylinder's axis sits at precisely the radius the specification asked for. In the
reference assembly the heat-exchanger cylinders are at $r=3.2000$ against a
specified \code{at\_radius} of $3.200$, and the pumps at $r=3.3690$ against
$3.369$. The simplified model is radially representative of the built geometry
by construction, not by adjustment.
\end{caveat}

\subsection{Overlap detection}

Equivalent cylinders are larger than the components they replace, so they can
overlap where the real components do not. \code{\_check\_overlaps()} tests every
pair for axis separation less than $r_1+r_2$ \emph{and} intersecting $z$ ranges,
and warns with the radial penetration and $z$ overlap, naming
\code{spec["hom\_cylinder"]} as the repair. The reference assembly reports no
overlaps.

% =====================================================================
\section{The assembly-level cells}
% =====================================================================

\subsection{Vessel}

The vessel is kept native rather than homogenised, because it \emph{is} the pool
boundary. It is modelled as a cylindrical annulus between $r_{\mathrm{in}}$ and
$r_{\mathrm{out}}$ spanning the pool height, plus a flat bottom disc whose
thickness is \emph{solved} so that the modelled steel equals the CAD steel
volume including the torispherical head:
\[
t_{\mathrm{bottom}} \;=\;
  \frac{V_{\mathrm{CAD}} - \pi\,(r_{\mathrm{out}}^2-r_{\mathrm{in}}^2)\,H}
       {\pi\,r_{\mathrm{out}}^2} .
\]
A negative result --- the annulus alone already exceeding the CAD steel --- is
clamped to zero and recorded in \code{qa.notes} rather than failing.

\subsection{Pool}

The pool is the vessel bore over the height from the lowest component cylinder
to the underside of the top plate, minus the volume every component cylinder
carves out of that range. Into that remainder the smeared components' steel is
dissolved, and the whole mix passes through the same
\code{homogenise\_volumes()} used for components --- so the pool's composition
conserves atoms on identical terms.

If a component cylinder extends radially past the bore, the carve-out is
approximated and a note is recorded. The reference assembly produces no such
note.

\subsection{Top plate}

The plate cell is the slab minus the penetrating cylinders; the CAD steel volume
is then smeared through it. Because the CAD volume is measured after the hole
pre-pass, the real penetrations are already excluded from it.

% =====================================================================
\section{The output model}
% =====================================================================

\code{homogenise\_assembly()} \emph{returns a Python dictionary}. It writes a
file only when \code{out\_path=} is given; the object is the same either way. It
contains geometry \emph{and} composition for every cell:

\begin{center}
\begin{tabular}{@{}lp{10.2cm}@{}}
\toprule
key & contents \\
\midrule
\code{components} &
  list of the equivalent cylinders: \code{\{x, y, z\_bottom, height, radius\}},
  \code{number\_densities}, and a per-component \code{qa} block holding the zone
  volumes, filler, unfilled remainder, per-material volumes, closure and worst
  atom error \\
\code{vessel} &
  \code{r\_inner}, \code{r\_outer}, \code{z\_bottom}, \code{z\_top},
  \code{bottom\_plate\_thickness}, \code{number\_densities}, and the CAD steel
  volume it conserves \\
\code{pool} &
  \code{radius}, \code{z\_bottom}, \code{z\_top}, \code{volume},
  \code{number\_densities} (with the smeared steel already mixed in), and the
  list of what was smeared \\
\code{plate} &
  \code{radius}, \code{z\_bottom}, \code{z\_top}, \code{number\_densities}, and
  both the CAD steel volume and the model cell volume \\
\code{settings} &
  \code{particles}, \code{batches}, \code{inactive}, and \code{source\_cylinder}
  set to the core's envelope \\
\code{qa} &
  \code{overlaps}, \code{worst\_rel\_atom\_error}, \code{worst\_closure},
  \code{notes} \\
\bottomrule
\end{tabular}
\end{center}

\noindent
Every cell carries its own \code{number\_densities} in atoms/b$\cdot$cm ---
components, pool, vessel and plate alike. Nothing further is needed to define
the materials of a transport calculation.

% =====================================================================
\section{Visualisation}
% =====================================================================

\code{build\_cad(model)} turns the model back into buildable geometry: one
cylinder primitive per component and a pipe for the vessel, via
\code{to\_cad\_specs()}, fed through \code{assemble\_objects}; then the pool and
plate as $z$-slabs with every component cylinder cut out of them. Because these
cells do not overlap one another, \emph{the CAD assembly's own overlap detection
acts as an independent check on the equivalent-cylinder layout}.

\code{to\_cad\_specs()} is called by \code{build\_cad} and is exposed publicly
because its output --- plain dictionaries carrying \code{number\_densities} ---
is a convenient shape for an exporter. It omits the pool and plate, so an
exporter is better served by reading the model dictionary directly.

The most informative view superimposes both models, as shown in the quick start:
the real geometry solid, the equivalent cylinders translucent over it, each its
own branch of the viewer tree. The overhangs quantified in
Section~\ref{sec:audit} are immediately visible in that view.

% =====================================================================
\section{Verification: what can fail, and what cannot}
\label{sec:verify}
% =====================================================================

This section matters more than any other for a reader deciding how much to trust
the output.

\subsection{The checks that are true by construction}

The per-nuclide atom balance computes

\begin{lstlisting}[style=py]
n_hom[nuc] += (vol / envelope_volume) * n     # divide by V_cyl
...
cell = nd * envelope_volume                   # multiply by V_cyl again
rel_error = abs(cell - src) / src
\end{lstlisting}

\noindent
It divides by $V_{\mathrm{cyl}}$, multiplies by $V_{\mathrm{cyl}}$, and checks
that the answer matches --- algebraically an identity. It can only ever measure
floating-point round-off, and duly reports $\sim\!10^{-16}$ whatever the model
contains. The volume closure is the same: the filler is \emph{defined} as
$V_{\mathrm{cyl}}-\sum_z V_z$, so the per-material volumes are constructed to
sum to $V_{\mathrm{cyl}}$.

\begin{caveat}{Neither the atom balance nor the closure can detect a wrong model}
Both would report $10^{-16}$ if every equivalent cylinder were placed in the
wrong location, or sized wrongly, or built from the wrong solid. They verify
arithmetic, not physics. The report prints this warning next to the figure.
\end{caveat}

\subsection{The check that can fail}

\code{check\_against\_cad(model, specs)} rebuilds the real CAD assembly from the
specifications and compares two independently produced sets of numbers: each
component's world bounding box and volume against its equivalent cylinder. It is
the only function in the subsystem that requires \code{specs} as well as
\code{model}, precisely because it does not trust the model's own provenance.

Cases in which it fires:

\begin{itemize}
  \item \textbf{A zone function's hard-coded datum drifts from its builder.}
    \code{ihx\_zones} declares \code{z\_bottom} $=-\,$\code{lower\_plenum\_dome\_radius},
    a hand-written statement about where \code{create\_ihx()} starts the solid.
    Change the builder's $z$ origin and the envelope silently detaches from the
    component.
  \item \textbf{The placement convention changes.} \code{\_world\_cylinder()}
    assumes rotation about the local origin followed by translation. If that
    order ever changes, every cylinder moves.
  \item \textbf{A zone function bakes a world position into its envelope},
    violating the local-frame contract, so the component is translated twice.
  \item \textbf{The resolver and the homogeniser diverge} on where a component
    lands.
\end{itemize}

\noindent
In every one of these the atom balance still reports $10^{-16}$.

\subsection{A gap in the present check}

\code{check\_against\_cad} compares \code{z\_bottom} only. It reports
\code{worst\_dz\_bottom} and nothing about the top of the component or its
radial extent. The consequence is documented in Section~\ref{sec:overhang}: two
component types are not enclosed by their own equivalent cylinders, and the
check reports a worst placement error of $1.0\times10^{-7}$ regardless.
Extending it to compare $z_{\mathrm{top}}$ and radius is the single highest-value
improvement available to this subsystem.

% =====================================================================
\section{Worked example: the ESFR-SMR reference assembly}
% =====================================================================

All figures below come from
\file{examples\_reactor/example\_final\_esfr\_smr\_with\_materials.py}, whose
thirteen components are those of the CAD reference assembly with material
assignments added. Lengths are in metres; number densities in
atoms/b$\cdot$cm.

\subsection{The resulting cells}

\begin{center}
\begin{tabular}{@{}lrrrrr@{}}
\toprule
cell & $x$ & $y$ & $r$ & $z_{\mathrm{bottom}}$ & $z_{\mathrm{top}}$ \\
\midrule
\code{ihx\_1} & 3.200 & 0.000 & 0.800 & 2.000 & 10.770 \\
\code{ihx\_2} & $-1.600$ & 2.771 & 0.800 & 2.000 & 10.770 \\
\code{ihx\_3} & $-1.600$ & $-2.771$ & 0.800 & 2.000 & 10.770 \\
\code{pump\_1} & 1.685 & 2.918 & 0.675 & $-0.385$ & 11.615 \\
\code{pump\_2} & $-3.369$ & 0.000 & 0.675 & $-0.385$ & 11.615 \\
\code{pump\_3} & 1.685 & $-2.918$ & 0.675 & $-0.385$ & 11.615 \\
\code{diagrid} & 0.000 & 0.000 & 2.415 & $-0.460$ & 0.590 \\
\code{core} & 0.000 & 0.000 & 1.800 & 0.590 & 4.500 \\
\code{strongback} & 0.000 & 0.000 & 3.030 & $-1.702$ & $-0.460$ \\
\code{above\_core\_structure} & 0.000 & 0.000 & 2.356 & 5.612 & 10.808 \\
\midrule
\code{vessel} & --- & --- & 4.505 & $-1.702$ & 9.000 \\
\multicolumn{6}{@{}l}{\hspace{1em}\footnotesize bore $4.455$, equivalent flat bottom $0.04235$, CAD steel $17.7625$} \\
\code{pool} & --- & --- & 4.455 & $-1.702$ & 9.000 \\
\multicolumn{6}{@{}l}{\hspace{1em}\footnotesize volume $430.831$, with the redan's $4.311$ of \code{ss316} smeared in} \\
\code{plate} & --- & --- & 5.000 & 9.000 & 9.500 \\
\multicolumn{6}{@{}l}{\hspace{1em}\footnotesize CAD steel $32.1767$ smeared through a model cell of $25.3908$} \\
\bottomrule
\end{tabular}
\end{center}

\subsection{One component in detail}

The heat exchanger is the richest case. Its envelope is
$V_{\mathrm{cyl}}=\pi\,(0.800)^2\,(8.770)=17.633$, and the zone model measures

\begin{center}
\begin{tabular}{@{}llrl@{}}
\toprule
zone & role & volume & material \\
\midrule
\code{structure}  & solid & 1.5221 & \code{ss316} \\
\code{tube\_side} & fluid & 5.3367 & \code{na\_secondary} \\
\code{shell\_side}& fluid & 9.6787 & \code{na\_primary} \\
\midrule
\code{\_filler}   & ---   & 1.0956 & \code{na\_primary} \\
\bottomrule
\end{tabular}
\end{center}

\noindent
The zones sum to $16.5375$; the filler closes the remaining $1.0956$ to the
envelope exactly. Primary sodium therefore totals $9.6787+1.0956=10.7743$. The
steel volume fraction is $1.5221/17.633 = 0.08633$, so for iron
\[
n_{\mathrm{Fe56}} = 0.08633 \times 5.90\times10^{-2}
                  = 5.093\times10^{-3}\ \text{atoms/b}\cdot\text{cm},
\]
which is what the model reports. The full composition of that cell is
$\{\mathrm{Fe56}\ 5.093\!\times\!10^{-3},\ \mathrm{Cr52}\ 1.381\!\times\!10^{-3},\
\mathrm{Ni58}\ 6.906\!\times\!10^{-4},\ \mathrm{Mo98}\ 1.036\!\times\!10^{-4},\
\mathrm{Na23}\ 2.002\!\times\!10^{-2}\}$.

\subsection{Quality figures}

\begin{center}
\begin{tabular}{@{}lll@{}}
\toprule
figure & value & interpretation \\
\midrule
worst relative atom error & $1.63\times10^{-16}$ & true by construction \\
worst volume closure      & $1.57\times10^{-16}$ & true by construction \\
overlapping cylinder pairs& none & meaningful \\
carve-out notes           & none & no cylinder exceeds the bore \\
worst placement error     & $1.0\times10^{-7}$ & meaningful, but $z_{\mathrm{bottom}}$ only \\
module self-test          & $0$ error, $0$ drift over a $10^3$ unit change & meaningful \\
\bottomrule
\end{tabular}
\end{center}

\subsection{Inventory: envelope against component}

The ratio of envelope volume to real component volume shows how far each
material is diluted. It is not an error --- the filler accounts for the
difference --- but it indicates how much of each cell is coolant.

\begin{center}
\begin{tabular}{@{}lrrr@{}}
\toprule
component & CAD volume & envelope volume & ratio \\
\midrule
\code{ihx} (each) & 1.522 & 17.63 & 11.6 \\
\code{primary\_pump} (each) & 2.223 & 17.18 & 7.7 \\
\code{diagrid} & 1.491 & 19.25 & 12.9 \\
\code{core} & 39.80 & 39.80 & 1.0 \\
\code{strongback} & 25.81 & 35.82 & 1.4 \\
\code{above\_core\_structure} & 10.69 & 90.58 & 8.5 \\
\bottomrule
\end{tabular}
\end{center}

\noindent
The core is the only cell that is exactly its own envelope: its zone model
declares the specified radius and height, and the solid is that cylinder, so no
filler is involved and the cell is $100\%$ fuel.

% =====================================================================
\section{Audit: known limitations}
\label{sec:audit}
% =====================================================================

\subsection{Three components are not enclosed by their equivalent cylinder}
\label{sec:overhang}

Envelopes are declared from nominal parameters rather than measured from the
built solid. Comparing each component's local CAD bounding box against its
declared envelope:

\begin{center}
\begin{tabular}{@{}lrrrrrr@{}}
\toprule
& \multicolumn{3}{c}{radius} & \multicolumn{3}{c}{$z_{\mathrm{top}}$ (local)} \\
\cmidrule(lr){2-4}\cmidrule(lr){5-7}
component & envelope & CAD & excess & envelope & CAD & excess \\
\midrule
\code{ihx} & 0.800 & 1.385 & $+0.585$ & 7.985 & 8.585 & $+0.600$ \\
\code{primary\_pump} & 0.675 & 1.375 & $+0.700$ & 12.000 & 12.000 & --- \\
\code{diagrid} & 2.415 & 2.434 & $+0.018$ & --- & --- & --- \\
all others & \multicolumn{6}{l}{enclosed exactly} \\
\bottomrule
\end{tabular}
\end{center}

\noindent
The heat exchanger's envelope is built from the plenum domes and bundle shell
only. The component also carries a riser rising \code{riser\_height} $=0.600$
above the upper dome, and a central pipe whose elbow swings out to $r\approx1.4$;
neither is in the envelope. The pump's nozzle elbow and flange overhang its
barrel radius, and this case is an explicit decision recorded in the source:
\emph{``Nozzles/flange overhang the barrel envelope: their atoms are still
counted, just smeared into the cylinder.''} The diagrid's excess is negligible.

Atoms remain conserved in every case --- the measured volume is the whole solid,
riser included --- but the material is \emph{relocated}, compressed into a
cylinder shorter and narrower than the real part. For an inventory or cost
calculation this is irrelevant. For a transport calculation it is a modelling
choice to be made knowingly: it is acceptable where the displacement is small
compared with a mean free path and the component is far from the tallies of
interest, and unacceptable where it is not.

\subsection{The top plate is over-carved}

This one differs in kind from the previous section, because volume is
reassigned \emph{between} cells rather than within one.

\begin{center}
\begin{tabular}{@{}lr@{}}
\toprule
quantity & value \\
\midrule
slab, $\pi r^2 t = \pi (5.0)^2 (0.5)$ & 39.270 \\
CAD steel after real penetrations & 32.177 \\
$\Rightarrow$ volume removed by the real penetrations & 7.093 \\
volume carved by the equivalent cylinders & 13.879 \\
$\Rightarrow$ excess carve & \textbf{6.786} \\
\bottomrule
\end{tabular}
\end{center}

\noindent
The equivalent cylinders remove nearly twice what the real penetrations do,
dominated by the above-core structure, whose cylinder is $r=2.356$ while its
neck through the plate is $r=1.109$. Two consequences follow:

\begin{itemize}
  \item The plate's steel is conserved but concentrated into a cell $27\%$
    smaller than the real steel volume, so the cell is \emph{denser} than the
    material it is made of: $n_{\mathrm{Fe56}}=7.604\times10^{-2}$ against
    \code{ss304}'s $6.00\times10^{-2}$.
  \item The excess $6.786$ of slab is filled by the components'
    \code{\_filler} --- sodium, above the pool, where the real machine has cover
    gas.
\end{itemize}

\noindent
Total steel atoms are correct. Sodium atoms are not, and the spurious sodium
sits directly above the core, where it is moderating material that does not
exist. The repair is to give the above-core structure a plate-region envelope
matching its neck, or to model the plate penetrations explicitly.

\subsection{The fuel composition does not match the geometry it is applied to}

\begin{caveat}{This dominates any $k$-eigenvalue result and must be fixed first}
\code{core\_smear} in \file{materials.py} was homogenised over a $95$~cm active
fuel column ($65$~cm fissile plus $30$~cm fertile, per the ESFR-SIMPLE D3.1
deliverable). The \code{CORE} specification is $3.910$~m tall, and its cell is
$100\%$ \code{core\_smear} with no filler. The model therefore faithfully
conserves approximately $3.910/0.95 = 4.1$ times the fissile inventory of the
real machine.
\end{caveat}

\noindent
This is the sharpest illustration of what conservation does and does not buy.
The transcription is exact: measured volume $39.80$, cylinder volume $39.80$,
per-nuclide error $10^{-16}$. Conservation guarantees the model contains what
was \emph{specified}; it cannot tell anyone the specification is right. The
repair is either to re-homogenise the composition over $3.910$~m, or to model
the core at its true active height --- a decision that depends on what the
$3.910$~m in the CAD represents. The same discrepancy dominates the cost
report, and is documented there as well.

\subsection{The model has no outer boundary}

The vessel stops at $z=9.0$ and the plate at $z=9.5$, but the heat-exchanger
cylinders reach $10.770$, the above-core structure $10.808$, and the pumps
$11.615$. Roughly $17$ of cell volume therefore lies outside any containing
cell. This is harmless in the JSON, and correct as a statement about the real
machine --- pumps and heat exchangers do penetrate the roof --- but a transport
code requires a bounded geometry. The exporter must wrap the model in a bounding
cylinder with a vacuum boundary condition and decide what fills the region above
the plate.

\subsection{Reconstructed fluid zones could drift silently}

Steel volumes are measured directly from the solid the assembly's own builder
produces, which makes them robust. The fluid zones --- \code{tube\_side},
\code{shell\_side}, \code{cavity} --- are \emph{reconstructions} from the same
parameters the builders use, and would drift if a builder changed without its
zone function following. The closure figure is the intended guard, since a
drifting reconstruction stops the zones summing to the envelope, but the guard
is weak: the filler absorbs any shortfall silently, and only a surplus large
enough to exceed the envelope raises an error.

\subsection{Tilted components cannot be represented}

A component with a non-zero roll or pitch has no $z$-aligned equivalent
cylinder. The homogeniser warns and places it upright. No component in the
reference assembly is tilted.

\subsection{No transport-code exporter exists}

See Section~\ref{sec:openmc}. The model is written and nothing consumes it yet.

\subsection{What was verified clean}

\begin{center}
\begin{tabular}{@{}lp{9.4cm}@{}}
\toprule
check & result \\
\midrule
module self-test & relative atom error $0$; number-density drift $0$ over a
  $10^3$ change of length unit \\
assembly atom balance & worst $1.63\times10^{-16}$ over all components \\
volume closure & worst $1.57\times10^{-16}$ \\
overlaps & none among the ten component cylinders \\
carve-out notes & none; no cylinder exceeds the vessel bore
  (largest outer reach $4.044$ against a bore of $4.455$) \\
placement & worst $z_{\mathrm{bottom}}$ deviation $1.0\times10^{-7}$ \\
\code{build\_cad} & builds all thirteen cells without error \\
call-order independence & \code{assemble\_objects} before
  \code{homogenise\_assembly} gives bit-identical plate and pool volumes \\
\bottomrule
\end{tabular}
\end{center}

% =====================================================================
\section{Handing the model to a transport code}
\label{sec:openmc}
% =====================================================================

The model dictionary is a neutral intermediate, deliberately not an OpenMC
input. Keeping the two apart is what allows the homogenisation to be computed
and viewed without a transport code installed at all. The translation belongs in
a separate module.

\subsection{The mapping}

\begin{center}
\begin{tabular}{@{}lp{9.4cm}@{}}
\toprule
model key & transport-code construct \\
\midrule
any \code{number\_densities} &
  a material: add each nuclide with atom fractions, then set the density as the
  summed atom density in atoms/b$\cdot$cm \\
\code{components[i].cylinder} &
  a $z$-cylinder intersected with two $z$-planes, forming one cell \\
\code{vessel} &
  the annulus between $r_{\mathrm{inner}}$ and $r_{\mathrm{outer}}$, plus the
  \code{bottom\_plate\_thickness} disc \\
\code{pool} &
  the bore slab minus every component region \\
\code{plate} &
  as the pool, over its own $z$ range \\
\code{settings} &
  batch and particle counts, and a cylindrical source distribution built from
  \code{source\_cylinder} \\
\bottomrule
\end{tabular}
\end{center}

\subsection{Three things the exporter must supply}

\begin{enumerate}
  \item \textbf{A length unit.} Transport codes generally work in centimetres,
    while these specifications are in metres: every length is multiplied by
    $100$. The number densities are \emph{not} rescaled --- they are already
    atoms/b$\cdot$cm, and Section~\ref{sec:units} showed them to be invariant
    under a change of the model's length unit. This is the point at which the
    unit-agnostic design finally commits, exactly as \code{build\_cad} does for
    STEP.
  \item \textbf{An outer boundary with a vacuum condition}, and a decision about
    what fills the space above the top plate. See the audit above.
  \item \textbf{Cross-section data and temperatures.} Nothing in the model
    records what temperature \code{na\_primary} at $0.83$~g/cm$^3$ corresponds
    to.
\end{enumerate}

% =====================================================================
\section{Extending the subsystem}
% =====================================================================

\paragraph{A new single-material component.} Nothing to do. Give it a
\code{"material"} and, if wanted, a \code{"filler"}; \code{generic\_zones()}
builds it through the same \code{build\_solid()} the assembler uses and gives it
its own bounding cylinder.

\paragraph{A component with internal structure.} Write a zone function returning
zones, envelope and optionally solids, and register it in
\code{MATERIAL\_ZONES}. Keep everything in the local frame. If the function
reconstructs a region rather than measuring it, keep it in step with the
builder, and treat the closure figure as the regression test.

\paragraph{A new material.} Add an entry to \code{MATERIALS} in either accepted
form. Specifications reference it by string; no import is required anywhere.

\paragraph{A component that should not become a cylinder.} Set
\code{"neutronics\_treatment": "smear"} to dissolve its material into the pool,
or \code{"skip"} to exclude it.

\paragraph{An envelope that overlaps a neighbour.} Set
\code{"hom\_cylinder": \{"radius": ...\}}. Atoms stay conserved; only the
dilution changes.

% =====================================================================
\section{Reproducing the results in this document}
\label{sec:repro}
% =====================================================================

\begin{lstlisting}
arch -arm64 .venv/bin/python homogenise.py     # atom-conservation self-test (no CadQuery)
\end{lstlisting}

\noindent
and, from a script holding the reference specification list:

\begin{lstlisting}[style=py]
from homogenise_assembly import (homogenise_assembly, check_against_cad,
                                 print_report, build_cad)

model = homogenise_assembly(SPECS, out_path="output/homogenised_model.json")
print_report(model, check_against_cad(model, SPECS))
build_cad(model)
\end{lstlisting}

\noindent
The specification list is that of
\file{examples\_reactor/example\_final\_esfr\_smr\_with\_materials.py}. Note that
importing that module executes it; to reuse only its specifications, read the
source and evaluate it up to the definition of \code{SPECS}.

% =====================================================================
\section{Summary}
% =====================================================================

The homogeniser converts the CAD specification list into a set of placed,
atom-conserving equivalent cylinders inside a sodium pool bounded by a natively
modelled vessel and top plate. Conservation is exact for every nuclide
independently, is unit-agnostic by construction, and is built on volumes
\emph{measured} from the CAD solids rather than estimated. Placement preserves
each component's radial position exactly. The output is a plain dictionary
carrying both geometry and composition for every cell, ready for a transport-code
exporter that has yet to be written.

Three cautions carry over into any use of the model. First, the atom balance and
the volume closure are true by construction and cannot detect a wrong model; only
\code{check\_against\_cad} can, and it currently compares $z_{\mathrm{bottom}}$
alone. Second, three components are not enclosed by their own equivalent
cylinders and the top plate is over-carved, so material is relocated by up to
$0.7$ in model units --- acceptable as a deliberate simplification, but a
decision rather than an accident. Third, and most importantly, the fuel
composition is applied to roughly four times the height it was derived for, and
no $k$-eigenvalue result should be believed until that is repaired.

\end{document}
